\section{Methodology}
\label{sec:method}

\subsection{Overview}
\label{sec:method:overview}

We design a four-condition experiment to measure how self-critique affects residual stream representations. For each math problem, we run a pipeline with four stages: (1)~generate an initial answer, (2)~generate a self-critique of that answer, (3)~generate a revised answer conditioned on the critique, and (4)~generate a control answer by re-prompting without critique. At each stage, we extract residual stream activations from all transformer layers, then analyze how representations shift between conditions.

\subsection{Model and Dataset}
\label{sec:method:model}

We use \qwen (7.6B parameters, 28 transformer layers, hidden dimension $d = 3{,}584$), loaded in float16 precision on a single NVIDIA RTX 3090. We select this model for its strong math reasoning capabilities and compatibility with activation extraction via forward hooks.

We evaluate on 200 problems from the \gsmk test split \cite{cobbe2021gsm8k}, a benchmark of multi-step arithmetic word problems with ground truth numerical answers. We use the first 200 problems in dataset order without selection bias.

\subsection{Pipeline Design}
\label{sec:method:pipeline}

Each problem is processed through four prompting conditions:

\para{Initial.} We prompt the model with: ``Solve the following math problem step by step. After your reasoning, provide your final answer as a number after `\#\#\#\# '.'' This establishes the baseline representation and answer.

\para{Critique.} We provide the model its own initial response and ask it to critique: ``You previously solved this problem: [question]. Your solution was: [initial\_answer]. Now carefully critique your own solution. Identify any errors.''

\para{Revised.} We prompt revision conditioned on the critique: ``Your original solution was: [initial\_answer]. Your self-critique identified these issues: [critique]. Now provide a revised solution.''

\para{Control.} We re-prompt the model with a slightly rephrased version of the initial prompt (``Please solve this carefully'') without any critique context. This isolates the effect of critique from simple prompt variation.

All conditions use greedy decoding (temperature $= 0$) with a maximum of 512 new tokens and random seed 42 for reproducibility.

\subsection{Activation Extraction}
\label{sec:method:activations}

We register forward hooks on each of the 28 transformer layers to capture residual stream hidden states. For each condition, we extract the activation vector $\vh_\ell \in \sR^{3584}$ at the \emph{last generated token} position for every layer $\ell \in \{0, 1, \ldots, 27\}$. This follows the probing methodology of \citet{marks2024geometry}, where last-token activations capture the model's final representation used for next-token prediction. The result is a tensor of activations with shape $(200, 4, 28, 3584)$---200 problems $\times$ 4 conditions $\times$ 28 layers $\times$ 3,584 dimensions.

\subsection{Analysis Methods}
\label{sec:method:analysis}

\para{Drift quantification.} We measure representational drift between conditions using cosine similarity $\cossim(\vh_\ell^{(i)}, \vh_\ell^{(j)})$ and L2 distance $\|\vh_\ell^{(i)} - \vh_\ell^{(j)}\|_2$ for each problem at each layer, where $i, j$ index conditions. We compare critique-induced drift (initial $\to$ revised) against control drift (initial $\to$ control) using paired Wilcoxon signed-rank tests with Benjamini-Hochberg FDR correction across layers.

\para{Linear probes.} We train $L_2$-regularized logistic regression classifiers ($C = 1.0$) to predict binary answer correctness from layer activations, using 5-fold stratified cross-validation. We report area under the ROC curve (\auc) for each condition and layer.

\para{PCA subspace analysis.} We apply PCA to activations from each condition at each layer, measuring variance explained by the top 3 principal components and the cosine alignment between principal components across conditions. High alignment indicates preserved subspace structure; low alignment indicates restructuring of the representational geometry.

\para{Drift--outcome correlation.} We compute point-biserial correlations between drift magnitude (cosine distance) and binary answer change (whether the revised answer differs from the initial answer) to test whether larger drift predicts behavioral change.

\para{Statistical testing.} We use McNemar's test for paired accuracy comparisons between conditions. All $p$-values are reported with Benjamini-Hochberg FDR correction where multiple comparisons are performed. We report Cohen's $d$ effect sizes throughout.
